% !TEX root = ../LosningsforslagTRES0312.tex
\chapter{Kinematikk}
\label{LF:Chapter:kinematikk}
Dette kapittelet inneholder løsningsforslag til oppgavene i Kapittel~\ref{komp-Chapter:kinematikk}, \textit{Kinematikk}, i kompendiet.

\oppgaver{Gjennomsnittsfart i rettlinjet bevegelse}{k3oppg1}{%
%
En bil kjører $\SI{120}{\kilo\metre}$ på $\SI{1.5}{\hour}$.
\begin{enumerate}[a)]
    \item Finn gjennomsnittsfarten i km/t.
    \item Gjør om til m/s.
    \item En annen bil kjører samme strekning på $\SI{90}{\minute}$. Hvem kjørte raskest?
\end{enumerate}
%
}

\svar[title={Løsningsforslag}]{%
\noindent
\textbf{a)} Vi bruker definisjonen av gjennomsnittsfart, $\bar v = \Delta s/\Delta t$ (se \textit{Gjennomsnittsfart og gjennomsnittshastighet} i Kapittel~\ref{komp-Chapter:kinematikk}):
\[
    \bar v = \frac{\SI{120}{\kilo\metre}}{\SI{1.5}{\hour}} = \doubleunderline{\SI{80}{\kilo\metre\per\hour}}
\]

\smallskip\noindent
\textbf{b)} Vi regner om med $\SI{1}{\kilo\metre\per\hour} = \dfrac{\SI{1000}{\metre}}{\SI{3600}{\second}}$:
\[
    \bar v = 80\;\si{\kilo\metre\per\hour}\cdot\frac{\SI{1000}{\metre}}{\SI{3600}{\second}} = \doubleunderline{\SI{22.2}{\metre\per\second}}
\]

\smallskip\noindent
\textbf{c)} Her må vi passe på enhetene: $\SI{90}{\minute} = \SI{1.5}{\hour}$, altså akkurat samme tid som den første bilen brukte! Den andre bilens gjennomsnittsfart blir dermed
\[
    \bar v_2 = \frac{\SI{120}{\kilo\metre}}{\SI{1.5}{\hour}} = \SI{80}{\kilo\metre\per\hour},
\]
som er \doubleunderline{nøyaktig like fort} som den første bilen. De to bilene kjørte altså like fort -- dette er en påminnelse om at man alltid bør sjekke enhetene nøye før man sammenlikner tall.
}

\oppgaver{Momentanhastighet}{k3oppg2}{%
%
Et legeme beveger seg etter posisjonslikningen $s(t) = 3t^2 + 2t$ (i meter og sekunder).
\begin{enumerate}[a)]
    \item For den som kan derivere: Vis at momentanfarten er gitt ved $v(t) = 6t + 2$.
    \item Hva er momentanhastigheten ved $t = \SI{2}{\second}$?
    \item Ved hvilken tid er momentanfarten $v = \SI{20}{\metre\per\second}$?
\end{enumerate}
%
}

\svar[title={Løsningsforslag}]{%
\noindent
\textbf{a)} Momentanfarten er den deriverte av posisjonen med hensyn til tid (se \textit{Momentanhastighet} i Kapittel~\ref{komp-Chapter:kinematikk}), $v(t) = ds/dt$. Vi deriverer $s(t) = 3t^2+2t$ ledd for ledd med potensregelen $\frac{d}{dt}(t^n) = nt^{n-1}$:
\[
    v(t) = \frac{d}{dt}\left(3t^2\right) + \frac{d}{dt}\left(2t\right) = 2\cdot 3t + 2 = \doubleunderline{6t+2} \quad\checkmark
\]

\smallskip\noindent
\textbf{b)} Vi setter inn $t=\SI{2}{\second}$:
\[
    v(2) = 6\cdot 2 + 2 = \doubleunderline{\SI{14}{\metre\per\second}}
\]

\smallskip\noindent
\textbf{c)} Vi løser likningen $v(t) = 20$ for $t$:
\[
    6t+2 = 20 \quad\Rightarrow\quad 6t = 18 \quad\Rightarrow\quad t = \doubleunderline{\SI{3}{\second}}
\]
}

\oppgaver{Konstant hastighet}{k3oppg3}{%
%
En syklist kjører med konstant fart $v = \SI{8.0}{\metre\per\second}$ i rettlinjet bevegelse.
\begin{enumerate}[a)]
    \item Hvor langt har syklisten kjørt etter $t = \SI{10}{\second}$?
    \item Hvor langt har syklisten kjørt etter $t = \SI{2}{\minute}$?
    \item Hvor lang tid tar det å tilbakelegge $s = \SI{2.0}{\kilo\metre}$?
\end{enumerate}
%
}

\svar[title={Løsningsforslag}]{%
\noindent
Farten er konstant, $v=\SI{8}{\metre\per\second}$, så vi bruker den enkle sammenhengen $s = v\cdot t$ (spesialtilfellet av posisjonslikningen i avsnitt~\ref{komp-sec:bevegelseslikninger} når $a=0$ og $s_0=0$).

\smallskip\noindent
\textbf{a)} Med $t=\SI{10}{\second}$:
\[
    s = \SI{8}{\metre\per\second}\cdot\SI{10}{\second} = \doubleunderline{\SI{80}{\metre}}
\]

\smallskip\noindent
\textbf{b)} Vi regner om $t=\SI{2}{\minute} = \SI{120}{\second}$:
\[
    s = \SI{8}{\metre\per\second}\cdot\SI{120}{\second} = \doubleunderline{\SI{960}{\metre}}
\]

\smallskip\noindent
\textbf{c)} Vi løser $s=v\cdot t$ for $t$:
\[
    t = \frac{s}{v} = \frac{\SI{2.0}{\kilo\metre}}{\SI{8}{\metre\per\second}} = \frac{\SI{2000}{\metre}}{\SI{8}{\metre\per\second}} = \doubleunderline{\SI{250}{\second}}
\]
}

\oppgaver{En oppgave til}{k3oppg4}{%
%
En båt beveger seg med fart $v=5.0$ knopp i Nord-Østlig retning. Én knop er definert som $1.852$ km/h.
\begin{enumerate}[a)]
    \item Finn farta i nordlig retning ($v_\text{N}$) og farta i østlig retning ($v_\text{Ø}$).
    \item Oppgi farten $v$ i km/h. 
\end{enumerate}
%
}

\svar[title={Løsningsforslag}]{%
\noindent
Vi regner først om farten til km/h, siden $1$ knop $=\SI{1.852}{\kilo\metre\per\hour}$:
\[
    v = 5.0\cdot\SI{1.852}{\kilo\metre\per\hour} = \SI{9.26}{\kilo\metre\per\hour}
\]

\smallskip\noindent
\textbf{a)} Nord-øst betyr en retning nøyaktig mellom nord og øst, altså en vinkel på $45^\circ$ fra nord-aksen. Vi dekomponerer farten (se \textit{Dekomponering av vektorer} i Kapittel~\ref{komp-Chapter:matematikk}) i nordlig og østlig komponent:
\[
    v_\text{N} = v\cos(45^\circ) = \SI{9.26}{\kilo\metre\per\hour}\cdot\cos(45^\circ) = \doubleunderline{\SI{6.55}{\kilo\metre\per\hour}}
\]
\[
    v_\text{Ø} = v\sin(45^\circ) = \SI{9.26}{\kilo\metre\per\hour}\cdot\sin(45^\circ) = \doubleunderline{\SI{6.55}{\kilo\metre\per\hour}}
\]
Komponentene blir like store, siden retningen ligger nøyaktig mellom nord og øst.

\smallskip\noindent
\textbf{b)} Farten i km/h fant vi allerede over:
\[
    v = \doubleunderline{\SI{9.26}{\kilo\metre\per\hour}}
\]
}

\oppgaver{Akselerasjon}{k3oppg5}{%
%
En bil øker farten fra $v_0 = \SI{20}{\metre\per\second}$ til $v = \SI{30}{\metre\per\second}$ på $\Delta t = \SI{5.0}{\second}$.
\begin{enumerate}[a)]
    \item Finn gjennomsnittsakselerasjonen.
    \item Er akselerasjonen positiv eller negativ? Hva betyr det?
    \item Dersom akselerasjonen er konstant, hva er farten etter ytterligere $\SI{3.0}{\second}$?
\end{enumerate}
%
}

\svar[title={Løsningsforslag}]{%
\noindent
\textbf{a)} Gjennomsnittsakselerasjonen er endring i fart delt på tidsintervallet (se \textit{Gjennomsnittsakselerasjon} i Kapittel~\ref{komp-Chapter:kinematikk}):
\[
    \bar a = \frac{\Delta v}{\Delta t} = \frac{v-v_0}{\Delta t} = \frac{\SI{30}{\metre\per\second}-\SI{20}{\metre\per\second}}{\SI{5}{\second}} = \doubleunderline{\SI{2.0}{\metre\per\second\squared}}
\]

\smallskip\noindent
\textbf{b)} Akselerasjonen er \doubleunderline{positiv}. Det betyr at farten øker -- bilen aksellererer (i vanlig forstand), den bremser ikke.

\smallskip\noindent
\textbf{c)} Dersom akselerasjonen er konstant, bruker vi hastighetslikningen (Likning 1 i avsnitt~\ref{komp-sec:bevegelseslikninger}), $v = v_0+at$, med ny starttid og startfart $v_0' = \SI{30}{\metre\per\second}$ og $t=\SI{3}{\second}$:
\[
    v = \SI{30}{\metre\per\second} + \SI{2.0}{\metre\per\second\squared}\cdot\SI{3}{\second} = \doubleunderline{\SI{36}{\metre\per\second}}
\]
}

\oppgaver{Sten i fritt fall}{k3oppg6}{%
%
\begin{enumerate}[a)]
    \item En sten faller utfor preikestolen. Hva er akselerasjonen til steinen (se bort ifra luftmotstand)
    \item En sten kastes rett oppover, og faller rett ned igjen. Hva er akselerasjonen til steinen (i) på vei opp, (ii) på toppen og (ii) på vei ned? Tegn figur og forklar kort med ord. 
\end{enumerate}
%
}

\svar[title={Løsningsforslag}]{%
\noindent
\textbf{a)} Når vi ser bort fra luftmotstand, er den eneste kraften som virker på steinen tyngdekraften. Alle legemer i fritt fall har da samme akselerasjon, tyngdens akselerasjon (se \textit{Fritt fall som eksempel} i Kapittel~\ref{komp-Chapter:kinematikk}):
\[
    a = \doubleunderline{\mathrm g = \SI{9.81}{\metre\per\second\squared}}\ \text{(rettet nedover)}
\]

\smallskip\noindent
\textbf{b)} Så lenge luftmotstand ses bort fra, er steinen i fritt fall under hele kastet -- både på vei opp, på toppen og på vei ned. Akselerasjonen er derfor \doubleunderline{konstant lik $\mathrm g$, rettet nedover, i alle tre tilfeller}. Det som endrer seg er \emph{farten}, ikke akselerasjonen:
\begin{center}
\begin{tikzpicture}[>=Stealth, scale=1]
    \node at (0,2.1) {\footnotesize\textbf{(i) På vei opp}};
    \draw[->, thick, blue] (0,0) -- (0,1.3) node[right, font=\scriptsize] {$\vec v$};
    \draw[->, thick, red] (0.7,0.6) -- (0.7,-0.6) node[right, font=\scriptsize] {$\vec a=\mathrm g$};

    \node at (3.2,2.1) {\footnotesize\textbf{(ii) På toppen}};
    \node[font=\scriptsize] at (3.2,0.65) {$\vec v = \vec 0$};
    \draw[->, thick, red] (3.9,0.6) -- (3.9,-0.6) node[right, font=\scriptsize] {$\vec a=\mathrm g$};

    \node at (6.4,2.1) {\footnotesize\textbf{(iii) På vei ned}};
    \draw[->, thick, blue] (6.4,0.6) -- (6.4,-0.7) node[right, font=\scriptsize] {$\vec v$};
    \draw[->, thick, red] (7.1,0.6) -- (7.1,-0.6) node[right, font=\scriptsize] {$\vec a=\mathrm g$};
\end{tikzpicture}
\end{center}
I alle tre tilfeller peker $\vec a$ nedover med samme størrelse $\mathrm g$ -- selv på toppen, der farten momentant er null, er akselerasjonen fortsatt $\mathrm g$ nedover (det er jo nettopp denne akselerasjonen som gjør at steinen begynner å falle igjen).
}

\oppgaver{Fritt fall}{k3oppg7}{%
%
En ball slippes fra ro fra en høyde $h = \SI{45}{\metre}$ over bakken. Husk at $\mathrm{g} = \SI{9.81}{\metre\per\second\squared}$.
\begin{enumerate}[a)]
    \item Finn tiden det tar før ballen treffer bakken.
    \item Finn farten rett før ballen treffer bakken.
\end{enumerate}
%
}

\svar[title={Løsningsforslag}]{%
\noindent
\textbf{a)} Ballen slippes fra ro, $v_0=0$, og faller med konstant akselerasjon $\mathrm g$. Vi bruker posisjonslikningen (Likning 2 i avsnitt~\ref{komp-sec:bevegelseslikninger}) med $s_0=0$ og $v_0=0$:
\[
    h = \tfrac{1}{2}\mathrm g t^2 \quad\Rightarrow\quad t = \sqrt{\frac{2h}{\mathrm g}} = \sqrt{\frac{2\cdot\SI{45}{\metre}}{\SI{9.81}{\metre\per\second\squared}}} = \doubleunderline{\SI{3.03}{\second}}
\]

\smallskip\noindent
\textbf{b)} Vi bruker hastighetslikningen $v=v_0+\mathrm g t$ med $v_0=0$:
\[
    v = \mathrm g\cdot t = \SI{9.81}{\metre\per\second\squared}\cdot\SI{3.03}{\second} = \doubleunderline{\SI{29.7}{\metre\per\second}}
\]
}

\oppgaver{Ballkast}{k3oppg8}{%
%
En ball kastes horisontalt fra en klippe $h = \SI{5.0}{\metre}$ over bakken med startfart $v_0 = \SI{10}{\metre\per\second}$. Bruk $\mathrm{g} = \SI{9.81}{\metre\per\second\squared}$.
\begin{enumerate}[a)]
    \item Finn tiden ballen bruker på å treffe bakken.
    \item Finn den horisontale avstanden fra klippen til der ballen lander.
    \item Finn farten (størrelsen av hastighetsvektoren) rett før ballen treffer bakken.
\end{enumerate}
%
}

\svar[title={Løsningsforslag}]{%
\noindent
Ballens bane er krumlinjet, så vi dekomponerer bevegelsen i en horisontal ($x$) og vertikal ($y$) retning, slik som i eksempelet \textit{Rallybil kjører utfor} (figur~\ref{komp-fig:rallybil}) i Kapittel~\ref{komp-Chapter:kinematikk}. Vi lar $x_0=0$ ved klippekanten og $y_0=h=\SI{5}{\metre}$ være høyden over bakken der ballen lander ($y=0$). Ballen kastes horisontalt, så
\[
    \vec v_0 = [v_{0x}, v_{0y}] = [\SI{10}{\metre\per\second},\,0], \qquad \vec a = [a_x,a_y] = [0,\,-\mathrm g].
\]
Komponentvis gir den generelle posisjonslikningen $\vec s(t) = \vec s_0+\vec v_0 t+\tfrac12\vec a t^2$:
\[
    x(t) = v_0 t, \qquad y(t) = h - \tfrac12 \mathrm g t^2.
\]

\smallskip\noindent
\textbf{a)} Ballen treffer bakken når $y(t)=0$:
\[
    0 = h - \tfrac12\mathrm g t^2 \quad\Rightarrow\quad t = \sqrt{\frac{2h}{\mathrm g}} = \sqrt{\frac{2\cdot\SI{5}{\metre}}{\SI{9.81}{\metre\per\second\squared}}} = \doubleunderline{\SI{1.01}{\second}}
\]

\smallskip\noindent
\textbf{b)} Den horisontale avstanden er $x(t) = v_0 t$:
\[
    x = \SI{10}{\metre\per\second}\cdot\SI{1.01}{\second} = \doubleunderline{\SI{10.1}{\metre}}
\]

\smallskip\noindent
\textbf{c)} Hastighetskomponentene idet ballen treffer bakken er $v_x=v_{0x}=\SI{10}{\metre\per\second}$ (uendret, ingen akselerasjon i $x$-retning) og
\[
    v_y = -\mathrm g t = -\SI{9.81}{\metre\per\second\squared}\cdot\SI{1.01}{\second} = \SI{-9.90}{\metre\per\second}.
\]
Farten (størrelsen av hastighetsvektoren) blir
\[
    v = \sqrt{v_x^2+v_y^2} = \sqrt{10^2+9.90^2} = \doubleunderline{\SI{14.1}{\metre\per\second}}
\]
}

\oppgaver{Prosjektilbevegelse}{k3oppg9}{%
%
Et prosjektil skytes ut fra bakken ($y_0 = 0$) med en kjent utgangshastighet $v_0$, i en vinkel $\theta$ over horisontalen. Se bort fra luftmotstand.
\begin{enumerate}[a)]
    \item Sett opp uttrykk for posisjonen til prosjektilet i $x$- og $y$-retning som funksjon av tiden $t$.
    \item Finn et uttrykk for tiden $t$ prosjektilet er i lufta, uttrykt ved $v_0$, $\theta$ og $\mathrm{g}$.

    \textit{Hint:} Bruk at prosjektilet returnerer til utgangshøyden når det lander, altså $y(t) = 0$.

    \item Bruk uttrykket for $t$ fra forrige punkt til å finne et uttrykk for hvor langt prosjektilet lander fra utgangspunktet (sluttposisjonen i $x$-retning).
    \item Skriv sluttposisjonen som en vektor $[x,\,y]$.
\end{enumerate}
%
}

\svar[title={Løsningsforslag}]{%
\noindent
Dette er samme type dekomponering som i eksempelet \textit{Straffespark} (figur~\ref{komp-fig:straffespark}) i Kapittel~\ref{komp-Chapter:kinematikk}, men nå helt generelt og symbolsk.

\smallskip\noindent
\textbf{a)} Prosjektilet starter i origo, $\vec s_0=[x_0,y_0]=[0,0]$, med utgangsfart dekomponert som $\vec v_0=[v_0\cos\theta,\,v_0\sin\theta]$, og akselerasjon $\vec a=[0,-\mathrm g]$ (kun tyngdekraften). Fra posisjonslikningen $\vec s(t)=\vec s_0+\vec v_0 t+\tfrac12\vec a t^2$, komponentvis:
\[
    x(t) = v_0\cos\theta\cdot t, \qquad y(t) = v_0\sin\theta\cdot t - \tfrac12\mathrm g t^2.
\]

\smallskip\noindent
\textbf{b)} Prosjektilet lander når $y(t)=0$:
\[
    0 = v_0\sin\theta\cdot t - \tfrac12\mathrm g t^2 = t\left(v_0\sin\theta - \tfrac12\mathrm g t\right)
\]
Dette gir enten $t=0$ (utgangspunktet) eller
\[
    \doubleunderline{t = \frac{2v_0\sin\theta}{\mathrm g}}
\]
som er tiden prosjektilet er i lufta.

\smallskip\noindent
\textbf{c)} Vi setter inn $t$ fra forrige punkt i uttrykket for $x(t)$:
\[
    x(t) = v_0\cos\theta\cdot\frac{2v_0\sin\theta}{\mathrm g} = \frac{2v_0^2\sin\theta\cos\theta}{\mathrm g}
\]
Ved å bruke den dobble vinkelen $\sin(2\theta)=2\sin\theta\cos\theta$ kan dette skrives kompakt som
\[
    \doubleunderline{x = \frac{v_0^2\sin(2\theta)}{\mathrm g}}
\]

\smallskip\noindent
\textbf{d)} Siden prosjektilet lander tilbake på utgangshøyden, $y=0$, blir sluttposisjonen
\[
    \doubleunderline{[x,\,y] = \left[\dfrac{v_0^2\sin(2\theta)}{\mathrm g},\ 0\right]}
\]
}

\oppgaver{Sykkel i sving}{k3oppg10}{%
%
En syklist kjører gjennom en sirkulær sving med radius $r = \SI{12}{\metre}$ og konstant fart $v = \SI{5.4}{\kilo\metre\per\hour}$.
\begin{enumerate}[a)]
    \item Gjør om farten til \si{\metre\per\second}.
    \item Finn sentripetalakselerasjonen til syklisten.
    \item Hva skjer med sentripetalakselerasjonen om syklisten dobler farten? Forklar svaret.
\end{enumerate}
%
}

\svar[title={Løsningsforslag}]{%
\noindent
\textbf{a)} Vi gjør om farten til SI-enheter:
\[
    v = \SI{5.4}{\kilo\metre\per\hour} = \doubleunderline{\SI{1.5}{\metre\per\second}}
\]

\smallskip\noindent
\textbf{b)} Vi bruker formelen for sentripetalakselerasjon (se \textit{Sentripetalakselerasjon}, figur~\ref{komp-fig:Sentripetalakselerasjon}, i Kapittel~\ref{komp-Chapter:kinematikk}):
\[
    a_\perp = \frac{v^2}{r} = \frac{(\SI{1.5}{\metre\per\second})^2}{\SI{12}{\metre}} = \doubleunderline{\SI{0.19}{\metre\per\second\squared}}
\]

\smallskip\noindent
\textbf{c)} Siden $a_\perp \propto v^2$, vil en dobling av farten ($v\to 2v$) gi
\[
    a_\perp' = \frac{(2v)^2}{r} = 4\cdot\frac{v^2}{r} = 4a_\perp,
\]
altså \doubleunderline{firedobles} sentripetalakselerasjonen -- den øker med kvadratet av farten, ikke proporsjonalt med farten.
}

\oppgaver{Sirkelbevegelse og omløpstid}{k3oppg11}{%
%
Månens gjennomsnittlige avstand fra jorda er $r = 3.84 \cdot 10^8\,\si{\metre}$, og omløpstiden er $T = 27.3\,\textrm{dager}$.
\begin{enumerate}[a)]
    \item Gjør om omløpstiden til sekunder.
    \item Finn sentripetalakselerasjonen månen opplever ved hjelp av formelen $a_\perp = 4\pi^2 r / T^2$.
    \item Beregn månens banehastighet $v = 2\pi r / T$ og verifiser svaret ved å bruke $a_\perp = v^2/r$ og sammenlikne med svaret du fikk i b).
\end{enumerate}
%
}

\svar[title={Løsningsforslag}]{%
\noindent
\textbf{a)} Vi gjør om omløpstiden til sekunder:
\[
    T = \SI{27.3}{\day}\cdot\frac{\SI{24}{\hour}}{\SI{1}{\day}}\cdot\frac{\SI{3600}{\second}}{\SI{1}{\hour}} = \doubleunderline{\SI{2.359e6}{\second}}
\]

\smallskip\noindent
\textbf{b)} Vi bruker formelen fra \textit{Omløpstid} (se Kapittel~\ref{komp-Chapter:kinematikk}):
\[
    a_\perp = \frac{4\pi^2 r}{T^2} = \frac{4\pi^2\cdot\SI{3.84e8}{\metre}}{(\SI{2.359e6}{\second})^2} = \doubleunderline{\SI{2.72e-3}{\metre\per\second\squared}}
\]

\smallskip\noindent
\textbf{c)} Banehastigheten er
\[
    v = \frac{2\pi r}{T} = \frac{2\pi\cdot\SI{3.84e8}{\metre}}{\SI{2.359e6}{\second}} = \doubleunderline{\SI{1.02e3}{\metre\per\second}} \approx \SI{1.02}{\kilo\metre\per\second}
\]
Som en verifikasjon bruker vi $a_\perp = v^2/r$:
\[
    a_\perp = \frac{(\SI{1.023e3}{\metre\per\second})^2}{\SI{3.84e8}{\metre}} \approx \SI{2.72e-3}{\metre\per\second\squared},
\]
som stemmer godt overens med svaret i b).
}

\oppgaver{Gjennomsnittsfart og gjennomsnittsakselerasjon}{k3oppg12}{%
%
En løper springer $\SI{400}{\metre}$ på $\SI{52}{\second}$.
\begin{enumerate}[a)]
    \item Hva er løperens gjennomsnittsfart, i $\si{\metre\per\second}$?
    \item Dersom løperen starter fra ro og oppnår toppfart $\SI{9.0}{\metre\per\second}$ etter de første $\SI{6.0}{\second}$, hva er gjennomsnittsakselerasjonen i denne perioden?
\end{enumerate}
%
}

\svar[title={Løsningsforslag}]{%
\noindent
\textbf{a)} Gjennomsnittsfarten er
\[
    \bar v = \frac{\Delta s}{\Delta t} = \frac{\SI{400}{\metre}}{\SI{52}{\second}} = \doubleunderline{\SI{7.69}{\metre\per\second}}
\]

\smallskip\noindent
\textbf{b)} Løperen starter fra ro ($v_0=0$) og når $v=\SI{9.0}{\metre\per\second}$ på $\Delta t=\SI{6.0}{\second}$. Gjennomsnittsakselerasjonen blir
\[
    \bar a = \frac{\Delta v}{\Delta t} = \frac{\SI{9.0}{\metre\per\second}-0}{\SI{6.0}{\second}} = \doubleunderline{\SI{1.5}{\metre\per\second\squared}}
\]
}

\oppgaver{Konstant akselerasjon -- oppbremsing}{k3oppg13}{%
%
En bil kjører med fart $v_0 = \SI{25}{\metre\per\second}$ og bremser med konstant akselerasjon $a = -\SI{4.0}{\metre\per\second\squared}$.
\begin{enumerate}[a)]
    \item Hvor lang tid tar det før bilen stopper?
    \item Hvor langt kjører bilen i løpet av oppbremsingen?
    \item Skisser fartsgrafen $v(t)$ for bilen fra den begynner å bremse til den stopper.
\end{enumerate}
%
}

\svar[title={Løsningsforslag}]{%
\noindent
\textbf{a)} Bilen stopper når $v=0$. Vi bruker hastighetslikningen (Likning 1 i avsnitt~\ref{komp-sec:bevegelseslikninger}), $v=v_0+at$, og løser for $t$:
\[
    0 = \SI{25}{\metre\per\second} + (-\SI{4.0}{\metre\per\second\squared})\cdot t \quad\Rightarrow\quad t = \frac{\SI{25}{\metre\per\second}}{\SI{4.0}{\metre\per\second\squared}} = \doubleunderline{\SI{6.25}{\second}}
\]

\smallskip\noindent
\textbf{b)} Vi bruker fart-strekning-likningen (Likning 3), $v^2=v_0^2+2a(s-s_0)$, med $v=0$ og $s_0=0$:
\[
    0 = (\SI{25}{\metre\per\second})^2 + 2\cdot(-\SI{4.0}{\metre\per\second\squared})\cdot s \quad\Rightarrow\quad s = \frac{625}{8.0} = \doubleunderline{\SI{78.1}{\metre}}
\]

\smallskip\noindent
\textbf{c)} Siden akselerasjonen er konstant, er $v(t)$ en rett linje som starter i $v_0=\SI{25}{\metre\per\second}$ og avtar lineært til $v=0$ ved $t=\SI{6.25}{\second}$:
\begin{center}
\begin{tikzpicture}[scale=1]
    \draw[->, thick] (-0.3,0) -- (5.0,0) node[right, font=\small] {$t$ [s]};
    \draw[->, thick] (0,-0.3) -- (0,3.2) node[above, font=\small] {$v$ [m/s]};
    \draw[thin] (0,2.5) -- (-2.5pt,2.5) node[left=4pt, font=\scriptsize] {$25$};
    \draw[thin] (4.5,2.5pt) -- (4.5,-2.5pt) node[below=3pt, font=\scriptsize] {$6.25$};
    \draw[headCol, thick] (0,2.5) -- (4.5,0);
    \node[below left=1pt, font=\scriptsize] at (0,0) {$0$};
\end{tikzpicture}
\end{center}
}

\oppgaver{Stein kastet fra et stup}{k3oppg14}{%
%
En stein kastes fra toppen av et stup med utgangsfart $v_0 = \SI{15}{\metre\per\second}$ i en vinkel $30^\circ$ over horisontalen. Stupet er $h = \SI{40}{\metre}$ over vannet. Se bort fra luftmotstand.
\begin{enumerate}[a)]
    \item Finn de horisontale og vertikale komponentene av $\vec{v}_0$.
    \item Finn hvor lang tid steinen er i luften før den treffer vannet.
    \item Finn hvor langt fra stupkanten steinen treffer vannet.
    \item Finn farten (størrelsen på hastigheten) steinen har i det den treffer vannet.
\end{enumerate}
%
}

\svar[title={Løsningsforslag}]{%
\noindent
Vi dekomponerer igjen bevegelsen, akkurat som i \textit{Rallybil kjører utfor} (figur~\ref{komp-fig:rallybil}). Vi lar $x_0=0$ og $y_0=h=\SI{40}{\metre}$ (høyden over vannet), med $y=0$ ved vannflaten.

\smallskip\noindent
\textbf{a)} Utgangsfarten dekomponeres med vinkelen $30^\circ$ over horisontalen:
\[
    v_{0x} = v_0\cos(30^\circ) = \SI{15}{\metre\per\second}\cdot\cos(30^\circ) = \doubleunderline{\SI{13.0}{\metre\per\second}}
\]
\[
    v_{0y} = v_0\sin(30^\circ) = \SI{15}{\metre\per\second}\cdot\sin(30^\circ) = \doubleunderline{\SI{7.5}{\metre\per\second}}
\]

\smallskip\noindent
\textbf{b)} Posisjonslikningen i $y$-retning er $y(t) = y_0+v_{0y}t-\tfrac12\mathrm g t^2$. Steinen treffer vannet når $y(t)=0$:
\[
    0 = \SI{40}{\metre} + \SI{7.5}{\metre\per\second}\cdot t - \tfrac12\cdot\SI{9.81}{\metre\per\second\squared}\cdot t^2
\]
Dette er en andregradslikning i $t$. Løser vi den med abc-formelen (og forkaster den negative løsningen, som ikke gir fysisk mening), får vi
\[
    \doubleunderline{t = \SI{3.72}{\second}}
\]

\smallskip\noindent
\textbf{c)} Horisontal avstand er $x(t) = v_{0x}\,t$:
\[
    x = \SI{13.0}{\metre\per\second}\cdot\SI{3.72}{\second} = \doubleunderline{\SI{48.3}{\metre}}
\]

\smallskip\noindent
\textbf{d)} Vertikal hastighetskomponent idet steinen treffer vannet er $v_y = v_{0y}-\mathrm g t$:
\[
    v_y = \SI{7.5}{\metre\per\second} - \SI{9.81}{\metre\per\second\squared}\cdot\SI{3.72}{\second} = \SI{-29.0}{\metre\per\second}
\]
Den horisontale komponenten er uendret, $v_x=v_{0x}=\SI{13.0}{\metre\per\second}$. Farten blir
\[
    v = \sqrt{v_x^2+v_y^2} = \sqrt{13.0^2+29.0^2} = \doubleunderline{\SI{31.8}{\metre\per\second}}
\]
Som en kontroll kan vi bruke energibetraktning (fritt fall-komponenten alene): $v=\sqrt{v_0^2+2\mathrm g h} = \sqrt{15^2+2\cdot9.81\cdot40}\approx\SI{31.8}{\metre\per\second}$, som stemmer.
}

\oppgaver{Sirkelbevegelse -- pariserhjul}{k3oppg15}{%
%
Et pariserhjul har radius $r = \SI{15}{\metre}$ og roterer med konstant fart slik at det tar $\SI{40}{\second}$ å fullføre én runde.
\begin{enumerate}[a)]
    \item Finn farten $v$ til en passasjer ved kanten av hjulet.
    \item Finn sentripetalakselerasjonen til passasjeren.
    \item Sammenlikn denne akselerasjonen med tyngdens akselerasjon $\mathrm{g}$. Er den stor eller liten i forhold?
\end{enumerate}
%
}

\svar[title={Løsningsforslag}]{%
\noindent
\textbf{a)} Vi bruker $v=2\pi r/T$ (se \textit{Omløpstid} i Kapittel~\ref{komp-Chapter:kinematikk}):
\[
    v = \frac{2\pi\cdot\SI{15}{\metre}}{\SI{40}{\second}} = \doubleunderline{\SI{2.36}{\metre\per\second}}
\]

\smallskip\noindent
\textbf{b)} Sentripetalakselerasjonen er
\[
    a_\perp = \frac{v^2}{r} = \frac{(\SI{2.36}{\metre\per\second})^2}{\SI{15}{\metre}} = \doubleunderline{\SI{0.370}{\metre\per\second\squared}}
\]

\smallskip\noindent
\textbf{c)} Forholdet til tyngdens akselerasjon er
\[
    \frac{a_\perp}{\mathrm g} = \frac{\SI{0.370}{\metre\per\second\squared}}{\SI{9.81}{\metre\per\second\squared}} \approx \doubleunderline{0.038},
\]
altså bare omtrent $3.8\,\%$ av $\mathrm g$. Sentripetalakselerasjonen er dermed \doubleunderline{liten} sammenliknet med tyngdens akselerasjon -- en passasjer i et pariserhjul kjenner knapt denne akselerasjonen.
}
